\documentclass[11pt]{article}
\usepackage[a4paper,margin=1in]{geometry}
\usepackage{amsmath,amssymb,amsthm}
\usepackage{booktabs}
\usepackage{hyperref}
\usepackage{microtype}
\usepackage[T1]{fontenc}
\usepackage{lmodern}

\newcommand{\RAKL}{RAKL}
\newcommand{\ImplementationSHA}{\texttt{UNBOUND}}
\newcommand{\SoftwareTests}{UNBOUND}
\newtheorem{definition}{Definition}
\newtheorem{proposition}{Proposition}
\newtheorem{theorem}{Theorem}
\newtheorem{corollary}{Corollary}
\newtheorem{lemma}{Lemma}

\title{Epistemic Mechanics: From Linguistic Claims to Evidence-Governed Scientific State}
\author{Sze Chun Yiu}
\date{10 August 2026}

\begin{document}
\setlength{\emergencystretch}{2em}
\maketitle

\begin{abstract}
Research agents built around large language models can make a claim easy to retrieve, repeat and act on long before the claim deserves scientific trust. We study this mismatch as a problem of state transition. In \RAKL{}, scientific objects live in an external typed state that records context, evidence, uncertainty, compatibility and scoped authority; proposal-generating models do not write to that state directly. Two additions address failures that arise even under this discipline. A transient workspace gives selected material broad computational access without granting either cross-context coherence or epistemic authority. Open-World Mechanism Discovery (OWMD) starts from unowned functions rather than framework vocabulary, forcing retrieval through function-only, historical, mathematical, implementation, citation, adversarial and freshness routes before a bounded closure claim is permitted. We prove the corresponding non-implications and the impossibility of certifying unrestricted open-world completeness from a finite search transcript. We also narrow the framework's mathematical terminology: the current atom--witness--path structure is a typed compatibility complex unless a scoped order and meet/join obligations are established separately. The contribution is therefore a formal and executable discipline for governing scientific state, not evidence of phenomenal consciousness, universal discovery or empirical scientific superiority.
\end{abstract}

\section{Introduction}

A research agent can place the same claim in every downstream context and still be wrong. It can also hold two individually plausible claims in active memory even when they cannot be reconciled under the same experimental conditions. These cases look similar if the system records only a scalar confidence score, yet they concern different questions: what is computationally available, what is coherent across context, and what the evidence licenses us to assert.

\RAKL{} was developed around that separation. Scientific objects are stored outside the language model in a typed, provenance-bearing state; model outputs enter as proposals and acquire authority only through explicit verification and promotion. That design solved one class of failure but exposed another. During development, recursive search repeatedly refined concepts inside the vocabulary already present in the framework and missed a functionally relevant global-workspace literature until that literature was introduced from outside. More recursion had not created a broader search space. It had deepened the same one.

This failure motivates the two mechanisms studied here. The first is a transient workspace for bounded selection and broadcast. It is deliberately weaker than a truth-maintenance system: selection changes which objects can influence computation, not which objects are scientifically licensed. The second is Open-World Mechanism Discovery (OWMD), which turns missing functions into vocabulary-independent search signatures and requires routes that can escape the framework's current terminology. Together they make a distinction that is easy to blur in agent architectures:
\[
\boxed{\text{computational access}\neq\text{cross-context coherence}\neq\text{epistemic authority}.}
\]

The paper also audits a word inherited from the implementation. The historical class `TypedKnowledgeLattice' stores typed atoms, compatibility witnesses and admissible constructive paths. Those objects are useful, but they do not by themselves supply an order-theoretic lattice. We therefore state the weaker structure explicitly and reserve lattice language for scopes in which the required order and meet/join laws are actually proved.

The results below are intentionally narrow. They concern transition permissions, compatibility witnesses and finite discovery protocols. They do not establish that an unrestricted literature has been exhausted, that workspace-like processing implies phenomenal consciousness, or that the present system improves scientific outcomes in real domains. Those questions require evidence of a different kind.

\section{Epistemic state and authority}

Let a local scientific chart be
\[
C_i=(U_i,\phi_i,\gamma_i,e_i,\alpha_i),
\]
where \(U_i\) is a scientific subdomain, \(\phi_i\) a representation, \(\gamma_i\) the relevant context and observation conditions, \(e_i\) evidence and provenance, and \(\alpha_i\) scoped scientific authority.

For a claim \(c\), authority is a coordinate tuple
\[
\alpha(c)=(G_c,R_c,M_c,I_c,D_c),
\]
covering grounding/provenance, representation/relation, mechanism ancestry, identification/bounding and decision or quantity-of-interest usability. The coordinates are certificate sets rather than a single confidence value. Under compatible scope, \(\alpha_1\preceq\alpha_2\) only when every relevant certificate coordinate in \(\alpha_1\) is contained in, or entailed by, the corresponding coordinate in \(\alpha_2\).

Canonical state is
\[
K_t=(\mathcal A_t,\mathcal T_t,\mathcal V_t,\mathcal E_t,\mathcal U_t,
\mathcal O_t,\mathcal F_t,\mathcal H_t^{-},\mathcal S_t,\mathcal G_t^K),
\]
containing atlas charts, typed transitions, surviving models, evidence and provenance, uncertainty objects, obstructions and residuals, research fibers, immutable negative history, discovery state and protected governance identities.

A model proposes rather than commits:
\[
G_\theta(K_t,a_t)\to\mathcal P_t.
\]
Verification then returns a typed outcome,
\[
\begin{aligned}
V(p,e,\gamma)\in\{&
\text{SUPPORTED},\ \text{REFUTED},\ \text{PARTIALLY IDENTIFIED},\\
&\text{BLOCKED},\ \text{CANNOT CHECK}\},
\end{aligned}
\]
and only a licensed update \(\mathcal U\) changes canonical state. The distinction is architectural rather than rhetorical: the proposer has no direct authority write path.

\section{Compatibility, gluing and the lattice boundary}

For overlapping charts, typed transitions
\[
T_{ij}^{\tau,\sigma}:\phi_i(U_i\cap U_j)\to\phi_j(U_i\cap U_j)
\]
carry a relation type \(\tau\), scope \(\sigma\), assumptions, evidence and error semantics. Contradiction is assessed only after the relevant context has been aligned:
\[
\operatorname{Contradict}_{\ell}(C_i,C_j)=
\operatorname{Overlap}(C_i,C_j)\land
\operatorname{Align}_{\ell}(C_i,C_j)\land
\neg\operatorname{Compatible}_{\ell}(C_i,C_j).
\]

The historical `TypedKnowledgeLattice' implementation does not add an order relation to these objects. Its operative content is more modest: typed atoms, pairwise compatibility witnesses and finite constructive paths.

\begin{definition}[Typed compatibility complex]
A typed compatibility complex is a triple
\[
\mathfrak C=(A,\kappa,\mathcal P),
\]
where \(A\) is a typed set of atoms, \(\kappa\) is a partial compatibility relation carrying explicit witnesses, and \(\mathcal P\) is a family of finite constructive paths whose declared pairwise obligations are witnessed.
\end{definition}

\begin{proposition}[Compatibility is not lattice structure]
The data of a typed compatibility complex is insufficient, by itself, to make \(\mathfrak C\) an order-theoretic lattice.
\end{proposition}

\begin{proof}
A lattice requires a partial order in which every relevant pair has a greatest lower bound and least upper bound. The data \((A,\kappa,\mathcal P)\) supplies neither such an order nor meet and join operations. Pairwise compatibility is also generally symmetric, whereas a non-trivial order is antisymmetric. Compatible atoms therefore need not possess a unique least upper bound. Any lattice claim requires additional structure.
\end{proof}

A local lattice interpretation remains possible when a scoped order \(\preceq_{\gamma,r}\) and its meet/join laws are established, or when a suitable closure operator is defined and its fixed points are shown to have the required lattice structure. The reference API consequently exposes `TypedCompatibilityComplex' as the semantically precise name while retaining the historical class for compatibility.

\section{A transient workspace for research cognition}

Persistent scientific state and active computational context solve different problems. At any step, only a small part of the archive can be useful to a bounded proposer, critic or search process. Let \(C_t\) be the candidates relevant to the current target and residual. A capacity-constrained gate selects
\[
W_t=G_{\kappa,\Pi_t^x}(C_t),\qquad |W_t|\le\kappa.
\]
The resulting frame records selected content \(S_t\), a typed broadcast map \(B_t\), a selection ledger \(L_t\), intervention handles \(X_t\), and lifetime/capacity state \(T_t\):
\[
W_t=(S_t,B_t,L_t,X_t,T_t).
\]
Downstream operators receive projections of \(W_t\) and return proposals. The default policy reserves capacity for challenge, novelty and negative-history material before filling the remaining slots by priority. This reservation is a simple guard against an attention mechanism that merely amplifies the system's current beliefs.

It is useful to name three kinds of ``global'' state separately:
\[
A_t(a)=\text{computational accessibility},\quad
C_t(a)=\text{atlas coherence},\quad
\alpha_t(a)=\text{epistemic authority}.
\]

\begin{proposition}[Authority-preservation invariant]
Suppose workspace transitions may write only workspace metadata and proposals, while increases in \(\alpha_t\) require the canonical verification/promotion path. Then a workspace-only transition cannot increase the authority of any scientific object.
\end{proposition}

\begin{proof}
The workspace transition has no write target in the authority coordinates and cannot create the certificate required by canonical promotion. Its projection onto canonical authority is therefore invariant. Hence \(\alpha_{t+1}(a)\not\succ\alpha_t(a)\) under a workspace-only transition.
\end{proof}

This is an invariant by construction, not evidence that a particular selection policy is scientifically optimal.

\begin{proposition}[Coactivation does not imply compatibility]
For arbitrary \(a,b\),
\[
a,b\in W_t\not\Rightarrow a\bowtie b.
\]
If a join operation is defined only for compatible objects in a scoped local lattice, then coactivation alone also does not establish that \(a\vee b\) exists.
\end{proposition}

\begin{proof}
Workspace admission depends on capacity, partition and computational priority rather than on an atlas compatibility witness. Two context-aligned but contradictory claims can therefore be selected simultaneously for adversarial comparison. Their coactivation leaves the compatibility predicate false and supplies no gluing witness.
\end{proof}

Three small counterexamples make the separation concrete. First, a critic may deliberately place two contradictory claims in the same workspace; access then increases while coherence fails. Second, two observational claims may glue cleanly across coordinates while still lacking the mechanism or identification evidence needed to increase authority. Third, a false claim may receive the highest workspace priority precisely because the system is trying to falsify it. Computational prominence is therefore not a proxy for truth.

\begin{proposition}[Conservative workspace metadata]
Let \(K_t^+=(K_t,W_t,\Pi_t^{\mathrm{cog}})\) extend canonical state with workspace and cognitive-provenance metadata. If no canonical update transition is executed, then any finite sequence of workspace selection, broadcast and intervention operations leaves \(\pi_K(K_t^+)=K_t\) unchanged.
\end{proposition}

\begin{proof}
Each workspace operation writes only to \(W\), cognitive provenance or proposal objects. Composition preserves that restriction. The canonical projection is therefore unchanged until a separately verified update is applied.
\end{proof}

\section{Cognitive influence is not evidential support}

Workspace interventions can be useful even when the selected item is wrong. For an item \(a\), consider
\[
do(a\notin W_t),\qquad do(a\mapsto a'),\qquad do(w_t(a)=\lambda).
\]
If such an intervention changes later proposals or decisions, it establishes that \(a\) was computationally load-bearing for that process. It does not establish that \(a\) was true, independently observed or causally responsible for the phenomenon under study.

The implementation therefore keeps two provenance types separate:
\[
\Pi^{\mathrm{cog}}\neq\Pi^{\mathrm{evid}}.
\]
A cognitive edge records which active item influenced which downstream proposal, optionally under an intervention. An evidential edge records evidence-to-claim lineage through a verification identity. There is no default coercion from the first graph to the second.

\section{Open-World Mechanism Discovery}

The motivating search failure suggested a more general diagnostic. Before asking whether a subsystem is well implemented, ask which functions it must own. For subsystem \(M\), let
\[
\mathcal F_{\mathrm{req}}(M)=\{f_1,\ldots,f_n\}.
\]
A function counts as owned only when the corresponding record names a mechanism and its scope, preconditions, postconditions, evidence, executable tests and failure semantics. A subsystem name alone is not an ownership argument.

For an unowned or contested function, OWMD constructs a signature
\[
\sigma(f)=(I,O,C,R,D,X),
\]
covering inputs, outputs, constraints, characteristic relations, dynamics or control behaviour, and intervention/failure signatures. Search then branches from this signature rather than from a single framework label. The mandatory ensemble includes exact and lexical variants, a function-only route with core labels masked, historical precursors, mathematical equivalents, implementation analogues, methodology-inspiration retrieval, backward and forward citation expansion, literature bridges, adversarial alternatives, cross-language routes when relevant, and a final freshness scan.

This design has familiar antecedents. The vocabulary problem shows how often people use different names for the same object or function \cite{furnas1987}. Literature-based discovery exploits useful connections between bodies of work that are not directly linked \cite{swanson1986}. Methodology Inspiration Retrieval makes a related distinction between topical similarity and methods that can transfer across research problems \cite{garikaparthi2025}. OWMD adopts the narrower engineering lesson: at least one search route should be able to succeed without repeating the vocabulary that defined the current ontology.

Retrieved mechanisms keep their source and route provenance. They are then classified as equivalent, subsumed, complementary, conflicting, novel residual or unresolved. Equivalent or subsuming prior work narrows the novelty claim; unresolved work remains visible rather than being forced into an existing category. A bounded discovery workspace additionally reserves slots for remote, challenging, historical and fresh candidates before ordinary relevance fill.

\begin{definition}[Bounded discovery closure]
For a declared route family \(\mathcal R_Q\), source universe \(\mathcal D\), budget \(B\) and cutoff \(t^\star\), a function is bounded-discovery-closed when: it has a complete owner or an explicit unresolved fiber; every mandatory route is complete or explicitly inapplicable; at least one route satisfies lexical independence; citation expansion satisfies its declared stability rule; the freshness scan reaches \(t^\star\); omission and nearest-work equivalence audits are complete; and unresolved candidates are preserved.
\end{definition}

The qualifier ``bounded'' matters. A finite search procedure can establish that a declared protocol has been completed; it cannot establish that nothing relevant exists outside the protocol.

\begin{theorem}[Unrestricted open-world completeness is not finitely certifiable]
Assume the concept and source universe is unrestricted and no complete membership oracle is available. A finite OWMD transcript cannot certify that no relevant undiscovered mechanism exists outside all returned results.
\end{theorem}

\begin{proof}
A finite run observes finitely many queries and returned objects. Consider a second world that produces the same transcript but also contains a relevant mechanism described only in a vocabulary or source never reached by those routes. The finite transcript is identical in the two worlds, so it cannot distinguish the existence of the omitted mechanism. Universal non-existence is therefore not identifiable from that transcript.
\end{proof}

Bounded execution does terminate under ordinary engineering assumptions. If the mandatory route set is finite, every route has a finite budget and terminates within it, and assimilation performs finite work per retrieved object, then the audit is a finite sum of finite computations. It ends either with `BOUNDED\_CLOSED' or with explicit open obligations. We treat this as a protocol property rather than as a substantive theorem about discovery.

\section{Workspace lineage and the J-space boundary}

The selection-and-broadcast idea is not new. Blackboard architectures coordinated heterogeneous knowledge sources through shared control structures \cite{hayesroth1985,nii1986}; Global Workspace and Global Neuronal Workspace accounts emphasize limited capacity and broad availability \cite{dehaene2001,mashour2020}. The Consciousness Prior similarly proposes selecting a small subset from a larger pool and broadcasting it \cite{bengio2017}, while Shared Global Workspace architectures let specialist modules compete for a bandwidth-limited communication channel \cite{goyal2021}. These traditions are prior art for workspace selection and broadcast, not components for which we claim novelty.

Recent work by Gurnee et al. uses the Jacobian lens to identify language-model representations with functional properties associated with a global workspace, including verbal report, directed modulation, internal reasoning, flexible generalization and selectivity \cite{gurnee2026}. Their result is useful here for two reasons, both narrower than an identity claim.

First, the study concerns a functional notion of access and explicitly does not take a position on phenomenal consciousness. Workspace-like computation is therefore not, on its own, evidence of subjective experience. Second, its J-space construction is geometric rather than order-theoretic. At fixed sparsity \(k\), J-space is generated by sparse non-negative combinations of J-lens vectors and has the geometry of a union of \(k\)-dimensional cones. That construction supplies neither a partial order nor meet/join laws. The study also leaves open the mechanism by which representations enter J-space. We use it as motivation for selective-access and intervention tests, not as a theorem about the \RAKL{} epistemic substrate.

\section{A hidden-name regression}

`GWT-OMISSION-01' turns the motivating miss into a software contract. The benchmark withholds the terms `global workspace', `consciousness', `J-space', `blackboard' and the relevant author names. Instead, it exposes a functional description: many background processes operate in parallel; a small set is admitted competitively under a capacity bound; active material is broadly reusable; contents can persist or be evicted; interventions on active material can change later decisions; and prominence does not establish truth.

A scored retrieval run passes only if a lexically independent route carries provenance to the registered family containing blackboard architectures, GWT/GNW, the Consciousness Prior, shared neural workspaces and the J-space study, without leaking the withheld names into the independent query. The present software tests the scoring, provenance and leakage rules. It does not yet constitute a prospective demonstration that the retriever will recover a genuinely unknown future concept; that remains an empirical evaluation.

\section{Why one scalar is not enough}

The same separation appears in the quantities used to describe a workspace. Let
\[
\mathbf w_t(a)=(o_t,b_t,p_t,i_t,e_t)
\]
denote occupancy, broadcast breadth, persistence, intervention influence and eviction sensitivity. These are computational coordinates. There is no default monotone map \(\mathbf w_t(a)\mapsto\alpha_t(a)\).

A simple counterexample is enough. Give a false but adversarially important claim maximal workspace priority so that critics repeatedly inspect it. At the same time, omit a well-established calibration fact because it is irrelevant to the immediate quantity of interest. The false claim can dominate the workspace coordinates while having low epistemic standing; the calibration fact can have high authority while having zero occupancy. Any scalarization that interprets computational prominence as authority will misorder at least one such pair.

\section{Implementation correspondence and scope}

The reference implementation exposes the formal contracts in \path{src/rakl/open_world_discovery.py}, \path{src/rakl/workspace.py} and \path{src/rakl/compatibility_complex.py}, with executable invariants in the corresponding test modules. Publication builds are generated from the same checkout that runs those tests.

For auditability, the exact source subject is printed separately rather than embedded in prose:
\begin{center}
\ImplementationSHA
\end{center}
The same checkout reports \SoftwareTests{} passing software tests. These tests establish reference-mechanics behaviour; they are not a substitute for validation on real scientific tasks.

Several empirical coordinates therefore remain open: prospective hidden-concept retrieval, matched same-model scientific workflows, fresh self-evolution transfer, substantially different real-domain case studies, domain-expert assessment and independent peer review. Internal reviewer roles used during development are explicitly not counted as independent review.

\section{Conclusion}

The difficult part of an agentic research system is not only generating candidate explanations. It is controlling what those explanations are allowed to change. The architecture studied here keeps three operations apart: making an object computationally available, establishing that it is coherent with other context-bound objects, and granting it scientific authority. A transient workspace handles the first without inheriting the permissions of the other two. OWMD addresses a complementary failure mode by forcing important searches to leave the vocabulary that produced the current ontology, while making the limits of finite search explicit.

The resulting discipline is deliberately conservative. Prior work narrows novelty rather than disappearing behind new names; unresolved mechanisms remain open fibers; negative history persists; and mathematical terminology is licensed only where its defining obligations are met. That conservatism is a feature of the scientific state, not a limitation to be hidden by fluent text.

\begin{thebibliography}{99}
\small
\bibitem{hayesroth1985} B. Hayes-Roth. A blackboard architecture for control. \emph{Artificial Intelligence} 26:251--321, 1985. doi:10.1016/0004-3702(85)90063-3.
\bibitem{nii1986} H. P. Nii. The Blackboard Model of Problem Solving and the Evolution of Blackboard Architectures. \emph{AI Magazine} 7(2):38--53, 1986.
\bibitem{dehaene2001} S. Dehaene and L. Naccache. Towards a cognitive neuroscience of consciousness: basic evidence and a workspace framework. \emph{Cognition} 79:1--37, 2001. PMID:11164022.
\bibitem{mashour2020} G. A. Mashour, P. Roelfsema, J.-P. Changeux, and S. Dehaene. Conscious Processing and the Global Neuronal Workspace Hypothesis. \emph{Neuron} 105:776--798, 2020. doi:10.1016/j.neuron.2020.01.026.
\bibitem{bengio2017} Y. Bengio. The Consciousness Prior. arXiv:1709.08568, 2017.
\bibitem{goyal2021} A. Goyal et al. Coordination Among Neural Modules Through a Shared Global Workspace. arXiv:2103.01197, 2021.
\bibitem{gurnee2026} W. Gurnee et al. Verbalizable Representations Form a Global Workspace in Language Models. \emph{Transformer Circuits Thread}, published 6 July 2026; arXiv:2607.15495, posted 16 July 2026.
\bibitem{furnas1987} G. W. Furnas, T. K. Landauer, L. M. Gomez, and S. T. Dumais. The Vocabulary Problem in Human-System Communication. \emph{Communications of the ACM} 30:964--971, 1987. doi:10.1145/32206.32212.
\bibitem{swanson1986} D. R. Swanson. Fish oil, Raynaud's syndrome, and undiscovered public knowledge. \emph{Perspectives in Biology and Medicine} 30:7--18, 1986. doi:10.1353/pbm.1986.0087.
\bibitem{garikaparthi2025} A. Garikaparthi, M. Patwardhan, A. S. Kanade, A. Hassan, L. Vig, and A. Cohan. MIR: Methodology Inspiration Retrieval for Scientific Research Problems. In \emph{Proceedings of ACL 2025}, 28614--28659. doi:10.18653/v1/2025.acl-long.1390.
\end{thebibliography}

\end{document}
